\section{The Aggregator in Practice}\label{sec:agg}

The paper treats the aggregator as harmless: it ``does not use their own private information'' and only
merges signatures~\cite{qiao2023}. That is true, and it is why aggregation does not weaken security. But
merging is a sum, $z=\sum z_i$. If one signature is invalid (a faulty sensor, a corrupted link or an
attacker) the aggregate fails verification, and the sum no longer says which component was wrong. One
bad sensor thereby denies service to every honest sensor in its batch. This is a known property of
aggregation, and locating bad signatures in a batch is a studied problem~\cite{pastuszak2000}; the
paper's system model simply leaves it out. We confirmed it on our implementation (one corrupted signature
in 32 makes the whole aggregate fail), and our contribution is a measurement of what handling it costs.

\subsection{Three aggregator policies}
\prov{Implementation}.
\begin{itemize}
\item \textbf{BLIND} (the paper's model): merge everything. A single bad signature loses the batch and the
      culprit cannot be found.
\item \textbf{PREVERIFY}: verify each signature before merging and drop the bad ones.
\item \textbf{RETAIN}: merge optimistically but keep the individual signatures. If the aggregate fails,
      split the batch in halves and verify each half as an aggregate, recursively. A clean half costs one
      aggregate verification whatever its size, so $f$ faults among $n$ cost on the order of $f\log_2 n$
      verifications. One bad signature among 32 was located with 11 verifications.
\end{itemize}
None of the policies uses secret material, so the paper's security argument is unaffected; a test
inspects the code to confirm this.

\subsection{Which policy, when}
\prov{Experiment}, \prov{Ours}. We simulated batches of $n=200$ signatures on P-256, each signature bad
independently with a given probability, 5 trials per rate, using measured verification costs.
\Cref{tab:agg} reports the total aggregator and cloud work per signature actually delivered.

\begin{table}[ht]
\centering\small
\caption{Cost per delivered signature (ms), P-256, batches of 200, 5 trials per row.}\label{tab:agg}
\begin{tabular}{rrrrrl}
\toprule
Fault rate & Bad per batch & BLIND & PREVERIFY & RETAIN & Cheaper of the two that recover \\
\midrule
0.0\% & 0.0 & 0.0566 & 0.1645 & 0.1163 & RETAIN \\
0.1\% & 0.0 & 0.0566 & 0.1598 & 0.1163 & RETAIN \\
0.5\% & 0.8 & 0.1414 & 0.1612 & 0.2384 & PREVERIFY \\
2.0\% & 3.4 & batch lost & 0.1626 & 0.4058 & PREVERIFY \\
5.0\% & 7.8 & batch lost & 0.1705 & 0.4781 & PREVERIFY \\
10.0\% & 18.4 & batch lost & 0.1784 & 0.5795 & PREVERIFY \\
\bottomrule
\end{tabular}

\smallskip
\footnotesize At the 0.1\% rate no trial drew a bad signature, so the first two rows are both clean batches.
\end{table}

BLIND is cheapest while every batch is clean, but from a 2\% fault rate it delivered nothing in any trial.
RETAIN beats PREVERIFY only on clean batches; from about one bad signature per batch of 200, PREVERIFY is
cheaper, and its cost stays nearly flat up to 10\%. The crossover is low because verifying a signature on
its own costs only $1.7$ times its share of an aggregate verification ($0.098$~ms against $0.057$~ms per
signature), not the order of magnitude one might expect. That makes pre-verification affordable.

\paragraph{Recommendation.} Where sensors almost never fail, RETAIN is cheapest. Where an attacker can
submit bad signatures, the attacker chooses the fault rate, so PREVERIFY is the safe default: it costs
about $2.9$ times a blind aggregation on a clean batch and never loses one. An aggregator too constrained
to verify can still limit the damage by sub-batching: under our measured cost model the best sub-batch
size falls from 25 signatures at a 0.1\% fault rate to 2 at 10\%. This is an availability result, not a weakness in the
cryptography.
